\section{Setting}\label{sec:setting}

\paragraph{Circuits.}
A circuit $C$ on inputs $x_1,\dots,x_n$ is a directed acyclic graph whose
leaves are the inputs and the constants $0$ and $1$, and whose internal nodes
(gates) compute $\wedge$, $\vee$ (any fan-in), $\neg$, $\oplus$ or a threshold
function $[\sum_j w_j y_j \ge t]$ of their inputs; one gate is the output. The
weights $w_j$ and the threshold $t$ are rationals given in binary (a
floating-point number stands for its exact binary value), and the comparison is
exact, with no tolerance (our implementation evaluates threshold gates the same
way). $|C|$ is the number of wires. For a gate $g$,
$\sv(g)\subseteq[n]$ is the set of input variables in its cone (its
\emph{support}; always the syntactic one). Pushing negations to the inputs gives
the negation normal form (NNF) of $C$: an $\{\wedge,\vee\}$-circuit over literals
and constants. For $\{\wedge,\vee,\neg\}$-circuits it has at most two copies of
every gate, one per polarity; a $\oplus$-gate of fan-in $k$ expands to $O(k)$ gates
per polarity, a threshold gate whose weights, scaled to integers, have polynomially
bounded total $W=\sum_j|w_j|$ expands by a sequential counter with $O(kW)$ wires (a
negative weight acts on the complemented input), and other gates stay opaque. An
$\wedge$-gate of an NNF is \emph{decomposable} if its children have pairwise
disjoint supports, and an $\vee$-gate is \emph{deterministic} if its children are
pairwise inconsistent; an NNF all of whose $\wedge$-gates are decomposable is a
DNNF, and a d-DNNF if moreover all its $\vee$-gates are deterministic
\cite{darwiche2002map}.

\paragraph{Distributions.}
$D$ is a product distribution on $\bits^n$ with rational marginals
$\Pr[x_i=1]=q_i$, again given in binary (a floating-point marginal stands for its
exact value). The \emph{input size} of $(C,D)$ is $|C|$ plus the total bit length
of the $q_i$ and of the weights and thresholds of $C$; running times that mention
only $|C|$ count arithmetic operations on such numbers. Product
structure is what makes ``disjoint supports'' mean ``independent''. A
distribution given by a sampler circuit $S:\bits^m\to\bits^n$ reduces to the
uniform case: $\Pr_{x\sim S(U_m)}[C(x)=1]=\Pr_{u\sim U_m}[(C\circ S)(u)=1]$.
(The reduction can destroy structure; we do not rely on it.) We write
\[
  p \;=\; p_D(C) \;=\; \Pr_{x\sim D}[C(x)=1]
\]
and think of $p$ as tiny: $2^{-50}$, $n^{-\omega(1)}$, or $2^{-\Omega(n)}$.
For a set or function $E$ we write $D(E)=\E_D[E]$.

\paragraph{What counts as an answer.}
Because $p$ can be arbitrarily small, additive error is useless; we ask for
relative error. To let a procedure decline, we let it output intervals.

\begin{definition}[sound interval estimation]\label{def:sound}
An \emph{estimation procedure} gets $(C,D,\eps,\delta)$ and outputs an
interval $[L,U]\subseteq[0,1]$ (possibly the trivial $[0,1]$, meaning ``no
verdict''). It is \emph{sound} if for every circuit $C$ and every $D$,
\[
\Pr\big[p_D(C)\notin[L,U]\big]\le\delta .
\]
It is \emph{complete} on a class $\mathcal C$ if for all $C\in\mathcal C$,
with probability at least $1-\delta$ the interval is tight:
$U\le(1+\eps)L$. It is \emph{efficient} if it runs in time polynomial in the
input size, $1/\eps$ and $\log(1/\delta)$.
\end{definition}

A sound procedure that is complete on $\mathcal C$ gives, on $\mathcal C$, a
$(1\pm\eps)$-approximation of $p$ with probability $1-2\delta$; outside
$\mathcal C$ its output is still a valid (if loose) interval. Soundness
everywhere is what makes the ``check a property, then sample'' procedures of
the problem statement safe to use on circuits we know nothing about.

\paragraph{The baseline.}
Monte Carlo draws $x^{(1)},\dots,x^{(N)}\sim D$ and reports the fraction of
ones. A relative-error interval needs
$N=\Theta\big(\log(1/\delta)/(p\eps^2)\big)$ samples, so its cost is
\[
\cost_{\MC}=\Theta\Big(\frac{|C|\log(1/\delta)}{p\,\eps^2}\Big),
\]
exponential in the \emph{surprise} $\log_2(1/p)$. The
Dagum--Karp--Luby--Ross stopping rule achieves this without knowing $p$ in
advance, and their algorithm $\mathcal{AA}$ is optimal up to a constant factor
among procedures that only see i.i.d.\ samples of a $[0,1]$-valued variable
\cite{dklr2000} (see \cite{chen2007inverse} for a repaired proof of the stopping-rule
guarantee); for $0/1$ samples of mean $p\ll1$ the optimal cost is
$\Theta(\log(1/\delta)/(p\eps^2))$. Beating sampling therefore requires looking at
the circuit.

\begin{definition}[beating sampling]
A procedure beats sampling on a class $\mathcal C$ if it is sound, efficient,
and complete on $\mathcal C$, and $\mathcal C$ contains circuits with
$p=n^{-\omega(1)}$ (for most of our classes, $p=2^{-\Omega(n)}$). Its cost is then
polynomial in the input size whatever $p$ is, whereas that of Monte Carlo grows like
$1/p$, exponentially in the surprise.
\end{definition}

\paragraph{Prover-assisted estimation.}
In the second setting an unbounded prover sends one string $\pi$ (an
\emph{argument}) to a randomized polynomial-time verifier $V$, which then
flips coins (for instance to sample) and outputs an interval $[L,U]$.

\begin{definition}[argument system]\label{def:argument}
$V$ is \emph{sound} if for every circuit $C$, every $D$ and every string
$\pi$, $\Pr_{\text{coins of }V}\big[p_D(C)\notin V(C,D,\pi,\eps,\delta)\big]\le\delta$.
It is \emph{complete} on $\mathcal C$ if for every $C\in\mathcal C$ there is a
$\pi$ of size polynomial in the input size such that with probability $1-\delta$
the output is tight. One-sided variants (lower-bound or upper-bound arguments) are defined
in the obvious way.
\end{definition}

This is the non-interactive, statistically sound (Merlin--Arthur) model. We
will see in \cref{sec:limits} that interaction (or computational soundness in the
random-oracle model) would trivialise the problem, while in this model upper bounds on $p$ cannot be
certified for all circuits unless $\coNP\subseteq\MA$.
